% ---------------------------------------------------------------------------
% Pistonball environment
% ---------------------------------------------------------------------------
\begingroup
% Line breaking for paragraphs with long, unbreakable identifiers (local to
% this chapter): allow looser lines instead of overfull ones.
\setlength{\emergencystretch}{3em}\tolerance=1000
\chapter{Pistonball}\label{chap:pistonball}

Pistonball is a cooperative physics game. A row of pistons forms the floor of a
walled arena and a ball is dropped near the right wall; every piston is an
agent that moves up or down, and together the pistons must roll the ball to
the left wall. Each piston observes the ball only when it is close to it, and
the reward of a piston depends on the ball's progress while it observes the
ball. The environment is used to study cooperation under local observability,
credit assignment with local rewards, scalability in the number of agents and
the effect of the observation radius $\kappa$ on learning. The physics is
simulated with \pkg{pymunk}; \pkg{pygame} is only needed for rendering and
keyboard control (extra \code{pistonball}).

The class is \code{env\_lib.pistonball\_env.PistonballEnv}; the package also
exports \code{ManualPolicy}, \code{PistonballRenderer} and
\code{PistonballLayout}.

% ---------------------------------------------------------------------------
\section{Arena and physics}\label{sec:pistonball-physics}

All coordinates are pygame screen pixels with $y$ growing downwards. With $n$
pistons the screen is $W = 160 + 40n$ pixels wide and 560 pixels high; the
walls are 80 pixels from each screen edge, so the arena spans
$x \in [80, W - 80]$ and $y \in [80, 480]$. The left wall is the goal.

\begin{itemize}
  \item \textbf{Pistons.} Piston $i$ is 40 pixels wide with its left edge at
        $x = 80 + 40i$. Its height $y_i$ (the $y$ coordinate of the piston
        head) moves by 4 pixels per unit action and is limited to
        $y_i \in [387, 451]$, a travel of 64 pixels around the middle position
        419. At reset every piston starts at $451 - d_i$ with $d_i$ drawn
        uniformly from $\{0, 4, \dots, 28\}$.
  \item \textbf{Ball.} A disc of radius 40 pixels with mass
        \code{ball\_mass}, friction \code{ball\_friction} and elasticity
        \code{ball\_elasticity}. It starts at
        $x = W - 150 + \delta_x$, $y = 370 + \delta_y$, where
        $\delta_x \in \{-30, \dots, 30\}$ and $\delta_y \in \{-15, \dots, 15\}$
        are uniform integers when \code{random\_drop=True} and zero otherwise,
        with an angular velocity drawn uniformly from $[-6\pi, 6\pi]$\,rad/s
        when \code{random\_rotate=True}.
  \item \textbf{Simulation.} Gravity is 750\,px/s$^2$ downwards, walls and
        pistons have friction 0.64, and one environment step is one physics
        step of $\Delta t = 1/20$\,s. The pistons are kinematic bodies that are
        moved to their new height before the physics step; the ball speed is
        capped at 1000\,px/s and its position is kept inside the screen.
\end{itemize}

% ---------------------------------------------------------------------------
\section{Spaces}\label{sec:pistonball-spaces}

\paragraph{Action.} One value per piston, stacked in an array of shape
\code{(n,)}:
\begin{itemize}
  \item continuous (\code{continuous=True}, default):
        \code{Box(-1, 1, (n,), float32)}; the value is clipped to $[-1, 1]$ and
        moves the piston by $4 a_i$ pixels, where $+1$ moves it up;
  \item discrete (\code{continuous=False}): \code{MultiDiscrete([3] * n)} with
        0 = down, 1 = stay and 2 = up (4 pixels).
\end{itemize}
Actions of the wrong shape, continuous actions with NaN or infinite values, and
discrete values outside $\{0, 1, 2\}$ raise \code{ValueError}.

\paragraph{Observation.} \code{float32} array of shape \code{(n, 7)} in an
unbounded \code{Box}. Row $i$ is the observation of piston $i$
(Table~\ref{tab:pistonball-obs}). The piston under the ball has index
\begin{equation}
  p(x) = \operatorname{clip}\bigl(\lfloor (x - 85)/40 \rfloor,\; 0,\; n - 1\bigr),
\end{equation}
and piston $i$ observes the ball when $|i - p(x)| \le \kappa$, where $x$ is the
ball centre and $\kappa$ = \code{kappa}. Columns 2--6 are zero for pistons that
do not observe the ball. The method \code{observable\_mask()} returns the
boolean mask of the observing pistons; its optional argument \code{ball\_x}
evaluates the mask for another ball position.

\begin{table}[htbp]
  \centering\small
  \caption{Pistonball observation (row $i$). $(x, y)$ is the ball centre,
    $(v_x, v_y)$ its velocity in px/s and $\dot\phi$ its angular velocity.}
  \label{tab:pistonball-obs}
  \begin{tabularx}{\textwidth}{@{}cl>{\raggedright\arraybackslash}X@{}}
    \toprule
    Column & Value & Meaning \\
    \midrule
    0 & $(y_i - 419)/32$ & piston height in $[-1, 1]$; $+1$ is the lowest position \\
    1 & $(5 + 40i)/(40n)$ & piston position along the arena in $[0, 1]$ \\
    2 & $(x - 80)/(40n)$ & ball $x$ position (observers only) \\
    3 & $(y - 80)/400$ & ball $y$ position (observers only) \\
    4 & $v_x / 15$ & ball $x$ velocity (observers only) \\
    5 & $v_y / 8$ & ball $y$ velocity (observers only) \\
    6 & $\dot\phi / 8$ & ball angular velocity (observers only) \\
    \bottomrule
  \end{tabularx}
\end{table}

% ---------------------------------------------------------------------------
\section{Reward and episode end}\label{sec:pistonball-reward}

Let $b_t = \lfloor x_t - 40 \rfloor$ be the $x$ coordinate of the ball's left
edge after step $t$ and $\Delta b = b_{t-1} - b_t$ its leftward displacement.
The ball term rewards leftward progress at half rate and penalises rightward
motion in full:
\begin{equation}
  g_t =
  \begin{cases}
    0.5\,\Delta b, & \Delta b > 0,\\
    \Delta b, & \text{otherwise.}
  \end{cases}
\end{equation}
The local reward of piston $i$ is
\begin{equation}
  \begin{split}
  r_i = \tau
      &+ g_t\, \mathbf{1}\bigl[i \in O(b_{t-1}) \cup O(b_t)\bigr]
       + \rho_T\, \mathbf{1}\bigl[\text{goal},\; i \in O(x_t)\bigr]
       + \rho_L\, \mathbf{1}\bigl[\text{goal},\; i = 0\bigr] \\
      &+ \mu\, \frac{|\Delta y_i|}{4}\, \mathbf{1}\bigl[|\Delta y_i| > \epsilon_\mu\bigr],
  \end{split}
  \label{eq:pistonball-reward}
\end{equation}
where $O(z) = \{ i : |i - p(z)| \le \kappa \}$ is the set of pistons within
$\kappa$ hops of the piston under position $z$ (the ball term uses the left-edge
positions $b$, the goal bonus the ball centre $x_t$), $\tau$ =
\code{time\_penalty}, $\rho_T$ = \code{termination\_reward}, $\rho_L$ =
\code{leftmost\_piston\_reward}, $\mu$ = \code{movement\_penalty},
$\epsilon_\mu$ = \code{movement\_penalty\_threshold} and $\Delta y_i$ the
movement of piston $i$ in pixels. The team reward returned by \code{step()} is
$R = \sum_i r_i$.

\paragraph{Goal and episode end.} With \code{terminated\_condition=True}
(default) the goal is reached when the ball's left edge will reach the left wall
within one step, $x_t - 40 + v_x \Delta t \le 80.5$; the episode then
terminates. \code{truncated} becomes \code{True} after \code{max\_cycles}
steps. For example, with \code{leftmost\_piston\_reward=1.0} and the default
termination reward, piston 0 receives $0.5 + 1.0 = 1.5$ on reaching the goal,
the other pistons that observe the ball receive 0.5, and the remaining pistons
receive no bonus.

% ---------------------------------------------------------------------------
\section{Parameters}\label{sec:pistonball-params}

\begingroup\small
\begin{longtable}{@{}>{\raggedright\arraybackslash}p{0.33\textwidth}>{\raggedright\arraybackslash}p{0.1\textwidth}>{\raggedright\arraybackslash}p{0.51\textwidth}@{}}
  \caption{Constructor parameters of \code{PistonballEnv}.}\label{tab:pistonball-params}\\
  \toprule
  Parameter & Default & Description \\
  \midrule
  \endfirsthead
  \toprule
  Parameter & Default & Description \\
  \midrule
  \endhead
  \bottomrule
  \endfoot
  \code{n\_pistons} & \code{20} & Number of pistons (agents), $\ge 1$. \\
  \code{time\_penalty} & \code{-0.1} & Reward $\tau$ added to every piston at every step. \\
  \code{continuous} & \code{True} & Continuous actions in $[-1, 1]$; \code{False} for \code{MultiDiscrete([3] * n)}. \\
  \code{random\_drop} & \code{True} & Randomise the initial ball position ($\pm 30$\,px horizontally, $\pm 15$\,px vertically). \\
  \code{random\_rotate} & \code{True} & Random initial angular velocity in $[-6\pi, 6\pi]$\,rad/s. \\
  \code{ball\_mass} & \code{0.75} & Ball mass ($> 0$). \\
  \code{ball\_friction} & \code{0.3} & Ball friction ($\ge 0$). \\
  \code{ball\_elasticity} & \code{1.5} & Ball elasticity ($\ge 0$). \\
  \code{max\_cycles} & \code{125} & Episode length; reaching it sets \code{truncated}. \\
  \code{render\_mode} & \code{None} & \code{None}, \code{"human"} (pygame window) or \code{"rgb\_array"}. \\
  \code{movement\_penalty} & \code{0.0} & Reward $\mu$ per 4 pixels of piston movement, typically negative; 0 disables it. \\
  \code{movement\_penalty\_threshold} & \code{0.01} & Movements of at most this many pixels are not penalised. \\
  \code{kappa} & \code{1} & Observation radius $\kappa$ in pistons ($\ge 0$). \\
  \code{terminated\_condition} & \code{True} & End the episode when the ball reaches the left wall. \\
  \code{leftmost\_piston\_reward} & \code{0.0} & Bonus $\rho_L$ for piston 0 when the goal is reached. \\
  \code{termination\_reward} & \code{0.5} & Bonus $\rho_T$ for every piston that observes the ball when the goal is reached. \\
\end{longtable}
\endgroup

\paragraph{Physics presets.} The ball parameters change the character of the
task considerably. A heavy, damped ball (\code{ball\_mass=2.0},
\code{ball\_friction=0.8}, \code{ball\_elasticity=0.5}) rolls slowly and rewards
coordinated ramps; a light, bouncy ball (\code{ball\_mass=0.5},
\code{ball\_friction=0.1}, \code{ball\_elasticity=2.0}) is harder to control.

\paragraph{Info dictionary.} \code{step()} returns \code{"local\_rewards"} and
\code{"agent\_rewards"} (both the \code{float64} local rewards $r_i$, shape
\code{(n,)}) and \code{"total\_reward"} (their sum, equal to the returned
reward). The info dictionary of \code{reset()} is empty. \code{step()} before
\code{reset()} raises \code{RuntimeError}.

% ---------------------------------------------------------------------------
\section{Keyboard control}\label{sec:pistonball-manual}

\code{ManualPolicy(env, agent\_id=0, show\_obs=False)} lets a person control one
piston at a time in a window. Create the environment with
\code{render\_mode="human"} so that the pygame window exists and receives the
key presses, and call the policy once per step: \code{W} or \code{Up} raises
the selected piston, \code{S} or \code{Down} lowers it (hold the key to keep
moving), \code{A}/\code{Left} and \code{D}/\code{Right} select the previous or
next piston, \code{Backspace} resets the environment, and \code{Esc} or closing
the window sets \code{policy.running} to \code{False}. All other pistons stay
in place.

\begin{pycode}
from env_lib.pistonball_env import ManualPolicy, PistonballEnv

env = PistonballEnv(n_pistons=8, render_mode="human")
policy = ManualPolicy(env)
obs, info = env.reset(seed=0)
while policy.running:
    obs, reward, terminated, truncated, info = env.step(policy(obs))
    if terminated or truncated:
        obs, info = env.reset()
env.close()
\end{pycode}

% ---------------------------------------------------------------------------
\section{Rendering}\label{sec:pistonball-render}

The renderer draws with pygame and needs no image assets
(Figure~\ref{fig:pistonball}). It shows a gradient background with shaded walls
and a highlighted goal strip on the left wall, the pistons as rounded
rectangles (pistons that observe the ball in the accent colour with a glow),
the ball with shading, a rotation marker and a fading motion trail, a
heads-up display with the step counter, episode return, ball $x$ velocity and
number of observing pistons, and a progress bar with the ball's distance to the
goal. Frames have the screen size, $(160 + 40n) \times 560$ pixels. The
\code{"rgb\_array"} mode never initialises the display and works on headless
machines; \code{"human"} mode opens a window capped at 20 frames per second.
Colours follow the active theme (Section~\ref{sec:envlib-rendering}).

\begin{figure}[htbp]
  \centering
  \includegraphics[width=0.62\textwidth]{pistonball.png}
  \caption{Pistonball with 12 pistons (light theme) at the moment the ball
    reaches the goal. The three pistons that observe the ball ($\kappa = 1$)
    are highlighted; the ramp formed by the pistons rolls the ball to the
    left wall.}
  \label{fig:pistonball}
\end{figure}

% ---------------------------------------------------------------------------
\section{Registered id}\label{sec:pistonball-ids}

\envid{Pistonball-v0} creates \code{PistonballEnv} with its defaults (20
pistons, continuous actions, $\kappa = 1$). Creating it requires \pkg{pymunk};
rendering also requires \pkg{pygame}.

% ---------------------------------------------------------------------------
\section{Usage}\label{sec:pistonball-usage}

The following heuristic reads the ball position from the observing pistons and
shapes the pistons into a ramp that falls towards the left: pistons right of
the ball move up, pistons left of it move down.

\begin{pycode}
import numpy as np
from env_lib.pistonball_env import PistonballEnv

env = PistonballEnv(n_pistons=12, kappa=1)
obs, info = env.reset(seed=0)                     # obs: (12, 7) float32
episode_return = 0.0
while True:
    action = np.zeros(env.n_pistons, dtype=np.float32)
    observers = np.any(obs[:, 2:] != 0.0, axis=1)
    if observers.any():
        ball_x = obs[observers, 2].mean()          # (x - 80) / (40 n)
        offset = (obs[:, 1] - ball_x) * env.n_pistons  # in piston widths
        action = np.clip(2.0 * offset, -1.0, 1.0).astype(np.float32)
    obs, reward, terminated, truncated, info = env.step(action)
    episode_return += reward
    if terminated or truncated:
        break
print(env.frames, terminated, episode_return)     # goal after 99 steps
env.close()
\end{pycode}

The example script evaluates a similar heuristic against random actions,
records episodes and offers keyboard control:

\begin{shellcode}
python examples/pistonball_demo.py --save renders/pistonball.gif
python examples/pistonball_demo.py --discrete --kappa 2 --movement-penalty -0.05
python examples/pistonball_demo.py --manual --n-pistons 8
\end{shellcode}

Over 20 episodes with the default configuration the heuristic of the script
reaches a mean return of $857.1 \pm 254.3$ and brings the ball to the goal in
95\,\% of the episodes, whereas random actions obtain $-290.5$ and never reach
the goal.

% ---------------------------------------------------------------------------
\section{Reproducibility and performance}\label{sec:pistonball-perf}

\code{reset(seed=...)} seeds \code{self.np\_random}, from which the initial
piston heights, the ball position and its angular velocity are drawn; the
physics itself is deterministic. The random draws are made in the same order as
before version 1.0, so seeded episodes start from the same states as in earlier
versions.

Measured on the development machine with 20 pistons, a step without rendering
takes 0.06\,ms (1.11\,ms before 1.0, when the screen was drawn at every step).
Rendering an \code{"rgb\_array"} frame takes about 3\,ms.

% ---------------------------------------------------------------------------
\section{Changes in 1.0}\label{sec:pistonball-changes}

See Appendix~\ref{app:migration} for migration advice.

\begin{itemize}
  \item Nothing is drawn unless rendering is requested; pygame is initialised
        lazily. Together with vectorised observations and rewards this makes
        steps 16 to 30 times faster.
  \item In discrete mode the declared action space is
        \code{MultiDiscrete([3] * n)}, one entry per piston. The per-piston
        action arrays accepted by \code{step()} are unchanged; the former
        single-integer \code{Discrete(3**n)} declaration was removed.
  \item Actions are validated (shape, range, NaN), and \code{step()} before
        \code{reset()} raises.
  \item New vector renderer (shaded pistons, highlighted $\kappa$-hop
        observers, ball trail, heads-up display, progress bar); the PNG sprites
        were removed. The former drawing methods (\code{draw()},
        \code{draw\_background()}, \code{draw\_pistons()},
        \code{enable\_render()}) are deprecated no-ops.
  \item Compatible with pymunk 6 and 7.
\end{itemize}

\endgroup
